\documentclass[11pt,a4paper]{article}
% Common preamble for the hint sheets (\input from each file)
\usepackage[utf8]{inputenc}
\usepackage[T1]{fontenc}
\usepackage{lmodern}
\usepackage{amsmath,amssymb,amsthm,mathtools}
\usepackage[a4paper,margin=2.3cm]{geometry}
\usepackage{enumitem}
\usepackage{xcolor}
\usepackage[hidelinks]{hyperref}
\renewcommand{\leq}{\leqslant}
\renewcommand{\le}{\leqslant}
\renewcommand{\geq}{\geqslant}
\renewcommand{\ge}{\geqslant}
\setlength{\parindent}{0pt}
\setlength{\parskip}{0.35em}
% \hintstitle{Sheet title}{Class work / Homework}
\newcommand{\hintstitle}[2]{\begin{center}{\LARGE\bfseries Hints: #1\par}\vspace{0.5em}{\large #2\par}\end{center}
\noindent\textit{The hints to each problem are cumulative. Read one, try again, and open the next only if you are still stuck.}\par\medskip\hrule\medskip}
\newcommand{\problem}[1]{\par\bigskip{\large\bfseries Problem #1}\par\nopagebreak\smallskip}
\newcommand{\subpart}[1]{\par\smallskip\textbf{(#1)}\par\nopagebreak}
\newcommand{\hint}[1]{\par\smallskip\noindent\textcolor{black!70}{\textbf{Hint #1.}}\enspace}
\newcommand{\note}[1]{\par\smallskip\noindent\textit{Note.} #1\par}
\providecommand{\lt}{<}
\providecommand{\gt}{>}
\providecommand{\tightlist}{\setlength{\itemsep}{0pt}\setlength{\parskip}{0pt}}
\emergencystretch=3em

\begin{document}
\hintstitle{Number theory -- 2}{Homework}

\problem{1}
\hint{1} If $r$ is a root of $x^3-3x+1$, find another root that is a polynomial in $r$. Try the substitution $x=t+t^{-1}$.
\hint{2} Check the identity $(r^2-2)^3-3(r^2-2)+1=(r^3-3r-1)(r^3-3r+1)$. So $r^2-2$ is again a root.
\hint{3} Uniqueness forces $r^2-2=r$, i.e.\ $r\in\{2,-1\}$, and then $p\mid 3$. Conversely, mod $3$ the polynomial is $(x+1)^3$, so $a=2$ is the only root. For $p=2$ there are no roots. Answer: $p=3$.

\problem{2}
\hint{1} Rewrite the condition in terms of $j=I(k)$: since $jk\equiv1$, $I(k+1)\equiv\frac{j}{j+1}=1-\frac1{j+1}$.
\hint{2} Put $m=j+1$. The condition becomes $p+1-I(m)<m-1$. So the count $D$ equals $\#\{m\in[2,p-1]: m+I(m)>p+2\}$.
\hint{3} The involution $m\mapsto p-m$ sends $m+I(m)=S$ to $2p-S$. Values $S\in\{p-1,p,p+1\}$ occur at most $6$ times ($m^2\equiv-1$, $m^2\mp m+1\equiv0$), and $S=p\pm2$ never occurs. Hence $2D\ge p-1-6$, which beats $p/4-1$ for $p>10$. Check $p=5,7,11$ directly.

\problem{3}
\hint{1} For $N=2^k$ put $D(n)=b_{n+N}-b_n$. From the recurrence, $D(n+1)=D(n)\,(1+2b_n+2b_{n+N})$. So all $D(n)$ have the same $2$-adic valuation as $D(0)=b_N$, and $b_{2N}-2b_N=D(N)-D(0)=b_N(P-1)$ with $P=\prod_{n<N}(1+2b_n+2b_{n+N})$.
\hint{2} Prove by induction on $k$: (i) $v_2(b_{2^k})=k+1$, and (ii) $Q_k=\prod_{n<2^k}(1+4b_n)\equiv1+2^{k+1}\pmod{2^{k+2}}$. Base: $b_2=4$, $Q_1=5$.
\hint{3} By (i), $b_{n+N}\equiv b_n\pmod{2^{k+1}}$, so $P\equiv Q_k\pmod{2^{k+2}}$ and $v_2(P-1)=k+1$. This gives $v_2(b_{2N}-2b_N)=2k+2$. Also $b_{2N}=b_N(P+1)$ with $P\equiv1\pmod 4$, which proves (i) for $k+1$. For (ii), $Q_{k+1}\equiv Q_k^2\pmod{2^{k+3}}$ because $4D(n)\equiv0\pmod{2^{k+3}}$, and $(1+2^{k+1}u)^2\equiv1+2^{k+2}$.

\problem{4}
\hint{1} Composite $n$: if $q^2\mid n$, then $q$ and $2q$ can only follow $1$, which is impossible. If $n$ is square-free with $s\ge2$ prime factors, every even $a_i$ must be followed by a quadratic residue mod $n$. Count: $\prod\frac{p_i+1}2<\frac{n-1}2$.
\hint{2} Prime $n=2^k+1$: let $g$ be a primitive root. Then $b$ can follow $a$ exactly when $v_2(a)\le v_2(\mathrm{ind}\,b)$, since the $a$-th powers are the $2^{v_2(a)}$-th powers. The two statistics $v_2(a)$ and $\mu(b)=v_2(\mathrm{ind}\,b)$ have the same distribution on $\{1,\dots,2^k\}$; for example, $a\mapsto g^a$ maps one level to the other.
\hint{3} Join each level $A_i=\{v_2=i\}$ to $B_i=\{\mu=i\}$ by any bijection. This gives a disjoint union of cycles. Merge two cycles by swapping the successors of two vertices of the same level. A cycle that cannot be merged contains all of $A_\lambda$ and $B_\lambda$, where $\lambda$ is its minimal level. But each $B_i$ with $i\le k-2$ contains an odd number ($B_i$ is closed under $x\mapsto n-x$), and $B_{k-1}=\{2^k\}\to1$. So that cycle contains an odd number, $\lambda=0$, and odd vertices merge freely.

\problem{5}
\hint{1} For a Mersenne prime $q=2^p-1$, the order of $2$ is $p$, since $2^j-1<q$ for $j<p$. For $F_k$ with $k\ge2$ see Problem~1 of the class work: $s=2^{2^{k-2}}(2^{2^{k-1}}-1)$ is a square root of $2$ and has order $2^{k+2}$.
\hint{2} Coprimality: $F_0F_1\cdots F_{m-1}=F_m-2$, so a common prime factor of $F_j$ and $F_m$ divides $2$. But the $F_j$ are odd. Infinitely many pairwise coprime numbers $>1$ give infinitely many primes.

\problem{6}
\hint{1} The primitive roots are exactly $g^k$ with $\gcd(k,p-1)=1$, and there are $\varphi(p-1)$ of them.
\hint{2} Product: $k\mapsto p-1-k$ pairs these exponents without fixed points when $p>3$. So the sum of the exponents is $\equiv0\pmod{p-1}$ and the product is $1$.
\hint{3} Sum: use M\"obius inversion on the condition $\gcd(k,p-1)=1$. You get $\sum_{d\mid p-1}\mu(d)\sum_{x^{(p-1)/d}=1}x$. The inner sum, over all $e$-th roots of unity, is $0$ unless $e=1$, so only $d=p-1$ survives.

\problem{7}
\hint{1} Show: if a prime $q\nmid n$ divides $\Phi_n(a)$, then $\mathrm{ord}_q(a)=n$. Since $x^n-1$ is separable mod $q$, $a$ is not a root of any $\Phi_d$ with $d<n$, $d\mid n$. Hence $q\equiv1\pmod n$.
\hint{2} Take $a=n\cdot N!$. Then $\Phi_n(a)\equiv\pm1$ modulo $n\,N!$, and $|\Phi_n(a)|>1$, so it has a prime factor $>N$ not dividing $n$. Why the method stops there: the argument only ever produces primes whose order data force $q\equiv1\pmod n$. (A ``Euclidean'' polynomial proof exists only for classes $a$ with $a^2\equiv1\pmod n$; in general one needs Dirichlet's theorem.)

\problem{8}
\hint{1} Small cases suggest primes. If $m$ is prime and $k=m-2n\ge1$, then $k\binom nk=n\binom{n-1}{k-1}$, so $n\mid k\binom nk$, and $\gcd(n,k)=\gcd(n,m)=1$.
\hint{2} Even $m=2n>2$: $\binom n0=1$ is not divisible by $n$. Odd composite $m=pq$ with $p$ the least prime factor, $q\ge p\ge3$: take $n=(m-p)/2\in[m/3,m/2]$, so $\binom{n}{m-2n}=\binom np$.
\hint{3} $\binom np=\frac np\binom{n-1}{p-1}$, and $n-1\equiv-1\pmod p$, so $\binom{n-1}{p-1}\equiv1\pmod p$ (Lucas). Hence $v_p\binom np=v_p(n)-1$ and $n\nmid\binom np$. Answer: $m$ prime.

\problem{9}
\hint{1} Factor: $n^4+n^2+1=(n^2-n+1)(n^2+n+1)=g(n-1)g(n)$ with $g(t)=t^2+t+1$. Let $a(t)$ be the largest prime factor of $g(t)$. Note $g(t^2)=g(t)g(t-1)$.
\hint{2} $\gcd(g(t),g(t+1))=1$, so $a(t)\ne a(t+1)$. The required equality reads $\max(a(n-1),a(n))=\max(a(n),a(n+1))$, which holds whenever $a(n)$ is a local maximum.
\hint{3} The sequence $a$ is not eventually increasing, because $a(t^2)=\max(a(t),a(t-1))$. It is not eventually decreasing either: it would become constant, contradicting $a(t)\ne a(t+1)$. Hence it has infinitely many local maxima.

\problem{10}
\hint{1} If $p\not\equiv1\pmod5$, every residue is a fifth power: take $a=0$. In general let $s_r=\#\{(a,b):a^2+b^5\equiv r\}$, so that $\sum_r s_r=p^2$, and bound the number $N=\sum_r s_r^2$ of solutions of $a^2+b^5\equiv c^2+d^5$.
\hint{2} For fixed $(b,d)$, $a^2-c^2\equiv d^5-b^5$ has $2p-1$ solutions if the right side is $0$ and $p-1$ otherwise. The pairs with $b^5\equiv d^5$ number at most $5p-4$. Hence $N\le p^3+4p^2-4p$.
\hint{3} Scaling $(a,b)\mapsto(\lambda^5a,\lambda^2b)$ maps $S_r$ to $S_{\lambda^{10}r}$. So if some $s_r=0$ with $r\ne0$, at least $(p-1)/10$ values of $r$ have $s_r=0$. Cauchy--Schwarz then gives $N\ge\frac{10p^4}{9p+1}$, which contradicts the bound for $p>38$.
\end{document}
